\section*{Apêndice A: Artefatos e reprodução do painel}

Os artefatos empíricos deste estudo partem da base de descoberta \path{Instrumentos/Codigos/Artefatos/selecao_spring_boot/runs/20260920_java/population.csv} (95.894 registros únicos) e são produzidos pelos \emph{scripts} de \path{pipeline_oficial_tcc/}, executados na ordem a seguir. As execuções históricas anteriores permanecem preservadas separadamente e não integram os resultados.

\begin{itemize}
    \item \path{01_coleta_descoberta/coleta_spring_boot.py} aplica os filtros de metadados e seleciona 52.419 repositórios Java ativos; \path{02_coortes_e_shas/identificar_shas_eras.py} aplica o recorte de criação e define os 785 candidatos.
    \item \path{04_extracao_ck/extrair_snapshots_ck_ia.py} clona cada candidato com o histórico completo da \emph{branch} principal, fixa o \emph{commit} de referência da coleta, seleciona o \emph{snapshot} de cada era, executa o CK e inspeciona os vestígios de IA até cada \emph{snapshot}. As extrações brutas por classe e por método, com código de produção e de testes discriminados, ficam na pasta da execução, e as medições consolidadas em \texttt{measurements\_785\_candidatos.csv}.
    \item \path{03_higienizacao_reporeaper/inspecionar_build_testes.py} apura C3 e C4 na árvore de arquivos do estado atual de cada clone, e \path{reporeaper_curation.py} aplica os quatro critérios (634 aprovados).
    \item \path{04_extracao_ck/executar_extracao_ck.py} aplica o filtro de completude longitudinal e materializa os 537 repositórios válidos em \texttt{repositorios\_validos.csv} e os 1.631 \emph{snapshots} em \texttt{measurements\_validos.csv} (20 pioneiros em E1--E4, com 80 \emph{snapshots}, e 517 principais em E2--E4, com 1.551 \emph{snapshots}).
    \item \path{05_deteccao_ia/consolidar_deteccao_ia.py} reúne, por repositório válido, o resultado da detecção de IA em cada era e os tipos de sinal observados em E4 (\texttt{resultado\_deteccao\_ia\_validos.csv}).
    \item \path{06_tabelas_artigo/gerar_tabelas_artigo.py} gera deterministicamente as tabelas LaTeX em \path{Artigo/Generated/} a partir dos CSVs validados ($n = 20$ em E1--E4; $n = 517$ em E2--E4; e $N = 537$ em E4, com $n_a = 88$ e $n_b = 449$).
\end{itemize}

Cada linha de medição registra identificador, nome completo, contagem de estrelas, grupo temporal, era, SHA avaliado, data do \emph{commit}, status de extração, total de classes, arquivos Java distintos, LOC físico (contado uma única vez por arquivo, sem duplicação por classes internas), as sete médias do CK e o resultado da detecção de IA. Falhas de clonagem, ausência de \emph{commit} na era, ausência de classes Java e extrações parciais permanecem registradas com o respectivo status, sem conversão em zero, e classes com LCOM* indefinido são omitidas apenas do denominador dessa métrica. Cada etapa grava um manifesto JSON com parâmetros, contagens e \emph{hashes} SHA-256 para verificação independente.
